\section{Conclusion}
\label{sec:conclusion}

We measured whether a delegated negotiation agent tells its principal what the principal's own instructions
cost. In a testbed where that cost is exact and the same for every agent, agents across \nModels\ models kept
the client's requirement, protected the confidential figure and reached deals above the client's fallback, and in
most episodes their reports gave the client no way to learn that the requirement had blocked a better deal. When
asked afterwards, most agents could name the blocked offer and its value, so the omission was not a failure to
track it.

The report prompt matters more than the briefing does. Rewording the requirement made no detectable difference
in a small ablation, and telling the agent its transcript would be read made it write more without disclosing
more. A standing instruction that the client relies on the account helped about half the models, and a client
question about turned-down terms helped all of them. For the design of oversight, this means that a delegated
agent's report is best treated as an answer to the questions it was asked, and that the interface, rather than
the user, should know which questions those are.

Two limits apply. The rates here describe how often the cost of an instruction is conveyed when such a cost
exists; the live-counterparty arm suggests that such situations arise much less often unaided. And the principal
in H5 is a language model, so the downstream cost is measured on a simulated reader. Whether human principals
would act on the omitted information is untested here, and the environment, transcripts and measures are
released so that the question can be put to people against exactly these episodes.
